\documentclass[11pt,a4paper]{article}
% Shared preamble for the three companion documents.  Deliberately minimal:
% only packages a bare TeX Live carries, since every figure is built from
% LaTeX `picture` primitives rather than from TikZ or an external plotter.
\usepackage[utf8]{inputenc}
\usepackage[T1]{fontenc}
\usepackage{amsmath,amssymb,amsthm}
\usepackage{booktabs}
\usepackage{algorithm}
\usepackage{algpseudocode}
\usepackage{xcolor}
\usepackage[margin=2.4cm]{geometry}
\usepackage[hidelinks]{hyperref}

\newcommand{\Z}{Z}
\newcommand{\Dset}{D}
\newcommand{\E}{E(P_{\Dset})}
\newcommand{\R}{\mathbb{R}}
\newcommand{\Zint}{\mathbb{Z}}
\newcommand{\fig}[1]{\input{figs/#1}}

% the legend shared by every criterion-space figure
\newcommand{\figlegend}{%
  \par\smallskip\noindent{\footnotesize
  \textcolor{green!45!black}{$\blacksquare$}~efficient point\quad
  \textcolor{black!45}{$\square$}~dominated point\quad
  \textcolor{blue!60!black}{$\square$}~the box being settled\quad
  \textcolor{red!35}{$\blacksquare$}~what the cut removes\quad
  \textcolor{red!65!black}{$\square$}~the children left\quad
  \textcolor{red!65!black}{$\circ$}~$x_B$\quad
  \textcolor{orange!85!black}{$\circ$}~the centre cut on\par}}

\newtheorem{proposition}{Proposition}
\newtheorem{remark}{Remark}


\newtheorem{theorem}{Theorem}
\newtheorem{corollary}{Corollary}
\newtheorem{definition}{Definition}

\title{\bfseries A metaheuristic--exact hybrid for the \emph{complete} efficient
set of a multi-objective integer program,\\ with a linear fractional decision
criterion}

\author{}
\date{}

\begin{document}
\maketitle

\begin{abstract}
\noindent
Let $\Dset = \{x \in \Zint^{n}_{+} : Ax \leqslant b\}$ be bounded and let
$\Z_1,\dots,\Z_p$ be criteria to maximise, each a ratio of affine forms with a
strictly positive denominator on $\Dset$.  We address two questions about the
efficient set $\E$: producing \emph{all} of it, and maximising a further linear
fractional criterion $\Phi$ over it.

Existing methods for the second question work in \emph{decision} space,
truncating $\Dset$ with Sylva--Crema cuts.  Each cut is a disjunction and
costs $p$ binary variables and $p+1$ big-$M$ rows, so every sub-problem is
strictly larger than the last.  We keep the unexplored part of the problem in
\emph{criterion} space instead, as a list of boxes: truncation becomes a box
subtraction, every sub-problem is $\Dset$ plus a handful of ordinary linear
rows, and the model never grows.

Three results follow.  First, a characterisation that makes the complete
efficient set cheap: \textbf{if $v$ is non-dominated then every $x \in \Dset$
with $\Z(x) = v$ is efficient}, so $\E$ is exactly the union of the slices of
the non-dominated set and \emph{no efficiency test is paid on any slice point}.
Second, a slice is $\Dset$ plus $p$ \emph{linear} equations, for fractional
criteria as much as for linear ones; solving those equations eliminates $r$
variables before enumeration begins and is worth a median $49\times$ over
scanning with them as pruning rows.  Third, a hybridisation that pays: a Pareto
local search supplies verified efficient points which are placed in the box
list without a probe.  Over 150 instances it reaches the whole of $\E$ at
$1.31\times$ the pure exact method, faster on 139, with identical answers
throughout.

The hybridisation only pays when the seeds are handed over in the order the
exact search's own probe would meet them; in the order the archive produces
them the same hybrid is \emph{slower}.  We also report, at the same length,
four changes that remove a quarter of the sub-problems and buy nothing, and a
single-program completeness certificate that is sound and $4.5\times$ slower,
because the negative results delimit where hybridisation can help at all.
\end{abstract}

\section{Problem and notation}

\begin{equation}
(P_{\Dset}) \qquad
\text{``}\max\text{''} \ \ \Z_k(x) = \frac{c_k^{\top}x + a_k}{d_k^{\top}x + b_k},
\quad k = 1,\dots,p, \qquad x \in \Dset ,
\end{equation}
with $\Dset = \{x \in \Zint^{n}_{+} : Ax \leqslant b\}$ non-empty and bounded
and every denominator $D_k(x) = d_k^{\top}x + b_k$ strictly positive on
$\Dset$.  A point $x \in \Dset$ is \emph{efficient} when no $y \in \Dset$
satisfies $\Z(y) \geqslant \Z(x)$ with at least one strict inequality; $\E$
denotes the set of efficient points and $F = \Z(\E)$ the non-dominated set.
Linear criteria are the degenerate case $d_k = 0$, $b_k = 1$.

Given a decision maker's criterion
$\Phi(x) = (U^{\top}x + \alpha)/(V^{\top}x + \beta)$,

\begin{equation}
(P_E) \qquad \max \ \Phi(x) \quad \text{s.t.} \quad x \in \E .
\end{equation}

\noindent
$\E$ is a union of faces of $\Dset$: non-convex, with no explicit description.

Throughout, the data of each criterion is scaled to integers (numerator and
denominator by the same factor, leaving the ratio untouched).  That makes the
quantity
\begin{equation}\label{eq:ek}
  e_k(x ; \bar x) \;=\; D_k(\bar x)\,D_k(x)\,\bigl(\Z_k(x) - \Z_k(\bar x)\bigr)
\end{equation}
\emph{linear in $x$ and integer-valued}, so that
$\Z_k(x) > \Z_k(\bar x) \iff e_k(x;\bar x) \geqslant 1$ exactly --- no minimal
step has to be estimated, which is what makes the transposition from linear to
fractional criteria faithful rather than approximate.

\begin{remark}\label{rem:notZ}
$\Phi$ is a function of $x$ and \textbf{not} of $\Z(x)$.  Two efficient points
with the same criterion vector can have different $\Phi$.  Every criterion-space
construction below has to account for this, and it is the reason the transfer
of criterion-space search from linear to fractional objectives is not
immediate.
\end{remark}

\section{Where the cost of the decision-space method comes from}
\label{sec:cost}

One iteration of the published approach maximises $\Phi$ over a truncated
region $\Dset_l$, tests the maximiser for efficiency, banks the best $\Phi$ on
the criterion slice about to be destroyed, and truncates:
$\text{delete } \{x : \Z(x) \leqslant \Z(\tilde x)\}$.

Keeping the complement means keeping the $x$ for which \emph{some} criterion
strictly improves on $\tilde x$.  That is a disjunction, not a system of linear
inequalities, and it is modelled with one binary per criterion:
\begin{equation}\label{eq:cut}
  e_k(x;\tilde x) \geqslant 1 \cdot y_k + M_k (1 - y_k), \quad k = 1,\dots,p,
  \qquad \sum_{k=1}^{p} y_k \geqslant 1 ,
\end{equation}
with $M_k$ a lower bound on $e_k$ over the region.  \textbf{Every cut therefore
adds $p$ binaries and $p+1$ rows}, and the late sub-problems cost far more than
the early ones.  This fixes the cost model used throughout:
\[
  \text{total cost} \ \approx\ (\text{number of sub-problems})
  \times (\text{cost of one, which grows}) .
\]

\section{The same truncation in criterion space}
\label{sec:box}

Keep the unexplored region as a list of \emph{boxes} instead.  Deleting
$\{\Z \leqslant \Z(a)\}$ from a box is a box subtraction, and it has a disjoint
$p$-way decomposition:
\begin{equation}\label{eq:split}
  \text{child } k: \quad e_k(x;a) \geqslant 1
  \ \text{ and } \ e_j(x;a) \leqslant 0 \ \text{ for } j < k .
\end{equation}
Child $k$ asks that $\Z_k$ strictly improve on $a$ while no earlier criterion
does.  The children are disjoint, they cover the box minus
$\{\Z \leqslant \Z(a)\}$, and --- this is the point --- each is $\Dset$ plus a
handful of ordinary linear rows.  \textbf{The model never grows.}  A box is
written with the integer-valued rows \eqref{eq:ek} rather than with numeric
bounds on $\Z_k$, which is what carries the construction to fractional
criteria unchanged.

Each box also carries, for nothing, an interval per criterion:
child $k$ raises $lo_k$ to $\Z_k(a)$ and lowers $hi_j$ to $\Z_j(a)$ for $j<k$.
We use it in Section~\ref{sec:wall}.

\section{The complete efficient set}
\label{sec:complete}

\subsection{Enumerating the non-dominated set}

The enumeration keeps a list of boxes covering the unexplored region.  It pops
a box, \emph{probes} it by maximising a fixed strictly positive direction over
$\Dset$ intersected with the box, repairs the maximiser to an efficient point
$a$, records $\Z(a)$, and splits the box around $a$ by \eqref{eq:split}.

\begin{proposition}\label{prop:front}
The enumeration terminates and records every non-dominated vector exactly once.
\end{proposition}

\begin{proof}
\emph{Termination.}  The probe returns a point $x$ of $\Dset$ inside the box,
and the efficient point $a$ it is repaired to dominates it, so
$\Z(x) \leqslant \Z(a)$ and the split around $a$ removes at least $\Z(x)$ from
that box.  The unexplored region shrinks on every pass, and the non-dominated
set of a bounded integer program is finite.

\emph{Completeness.}  A non-dominated $v$ leaves the region only through a
split around some centre $a$ with $v \leqslant \Z(a)$.  If $v \neq \Z(a)$ that
says $a$ dominates $v$, which no non-dominated vector admits.  So $v$ leaves
only at the moment it is recorded; the loop ends only when the region is empty,
so every $v$ is recorded, exactly once since centres are de-duplicated by
criterion vector.
\end{proof}

\subsection{From the non-dominated set to the efficient set}

Completeness of $F$ is not completeness of $\E$: one vector can be attained by
several efficient points and the enumeration keeps one per vector.  The gap
closes on one observation.

\begin{theorem}\label{thm:slice}
If $v$ is non-dominated, then \textbf{every} $x \in \Dset$ with $\Z(x) = v$ is
efficient.
\end{theorem}

\begin{proof}
Suppose some $y \in \Dset$ dominated such an $x$.  Then
$\Z(y) \geqslant \Z(x) = v$ with at least one strict inequality, which says $v$
is dominated --- against the hypothesis.
\end{proof}

\begin{corollary}\label{cor:union}
$\displaystyle \E \;=\; \bigcup_{v \in F} \{\, x \in \Dset : \Z(x) = v \,\}$,
and \textbf{no efficiency test is paid on any slice point}.
\end{corollary}

\noindent
The dearest operation in the whole method --- the efficiency test --- is simply
not invoked on the slices.  What remains is to enumerate each slice.

\subsection{\texorpdfstring{A slice is $\Dset$ plus $p$ linear equations}{A slice is D plus p linear equations}}
\label{sec:elim}

Fixing $\Z(x) = v$ is linear in $x$ \emph{even for fractional criteria}: since
$D_k(x) > 0$ on $\Dset$, the equation
$(c_k^{\top}x + a_k)/(d_k^{\top}x + b_k) = v_k$ clears to
\begin{equation}\label{eq:slice}
  (c_k - v_k d_k)^{\top} x \;=\; v_k b_k - a_k ,
\end{equation}
with no sign condition and no case split.  A slice therefore carries no binary,
no big-$M$ and nothing that grows from one slice to the next.

These are equations, so \emph{solve} them rather than use them to prune.
Gaussian elimination in exact rationals expresses $r$ of the variables as
linear functions of the rest; substituting into every row of the model leaves
an equivalent system in $n - r$ unknowns, on which the enumeration runs.  Two
rows per eliminated variable keep it inside its own bounds.  At each leaf the
eliminated variables are rebuilt and the one condition that is \emph{not}
linear --- that they come out integer --- is checked there.  An inconsistent
system settles the slice as empty with no enumeration at all.  The search space
$(ub+1)^{n}$ becomes $(ub+1)^{n-r}$.

\begin{table}[h]
\centering
\caption{Pruning on the slice equations against solving them.  Fourteen
instances, $n = 7 \dots 12$; identical point sets throughout.}
\label{tab:elim}
\begin{tabular}{lrrr}
\toprule
instance & prune & eliminate & ratio \\
\midrule
$n=9$, $p=3$  & $29.71$\,s  & $0.31$\,s & $96.7\times$ \\
$n=10$, $p=5$ & $119.47$\,s & $0.92$\,s & $129.9\times$ \\
$n=12$, $p=3$ & $796.84$\,s & $8.34$\,s & $95.5\times$ \\
\midrule
\textbf{median of all fourteen} & & & $\mathbf{49.2\times}$ \\
\bottomrule
\end{tabular}
\end{table}

\noindent
Every coefficient is a rational and the elimination is exact, so this changes
the order of the enumeration and nothing about its output.

\subsection{\texorpdfstring{Pairing each vector with a maximiser of $\Phi$}{Pairing each vector with a maximiser of Phi}}

By Remark~\ref{rem:notZ} the front alone does not answer $(P_E)$.  Each vector
is therefore paired with the maximiser of $\Phi$ \emph{on its slice}, at one
extra integer program per vector.  That pairing has a consequence worth
stating: \textbf{the largest value on the front is the optimum of $(P_E)$}, so
the enumeration is a second, independent exact method for $(P_E)$ and serves as
a cross-check on the first.

\section{The hybridisation}
\label{sec:hybrid}

\subsection{What a seed can feed here, and what it cannot}

A seed for an \emph{optimisation} search feeds one mechanism: the bound.  A
search that already knows a good value prunes boxes whose bound cannot beat it.
\textbf{The enumeration has no bound and no incumbent} --- it must exhaust every
box --- so a seed of that kind is worth exactly nothing.

It has a different mechanism.  Every box costs a probe, an efficiency test and
sometimes a repair, all spent to find \emph{a vector}; a vector already known
needs none of them.  Seeding here is not ``start from a better value'' but
\emph{do not pay to discover what you already know}.

\begin{proposition}\label{prop:seedsound}
Let $S \subseteq \E$ be a set of efficient points.  Placing each $a \in S$ in
the box list by locating the unique box containing it and splitting that box
around $a$ by \eqref{eq:split}, then running the enumeration, returns the same
$F$ and the same $\E$, still complete by Proposition~\ref{prop:front}.
\end{proposition}

\begin{proof}
The split \eqref{eq:split} is disjoint and covers the box minus
$\{\Z \leqslant \Z(a)\}$, so a point whose vector has not been removed lies in
exactly one child; the location is therefore well defined, and a seed whose
vector has already been removed is dominated by an earlier one and is
discarded.  The operation performed on the list is exactly the one the
enumeration performs when its probe returns $a$: same centre, same rows, same
children.  The proof of Proposition~\ref{prop:front} appeals only to the
centres being feasible and to the region shrinking, both of which hold
verbatim.  Efficiency of the seeds is what makes the \emph{recorded} vectors
non-dominated, and is required for that reason alone.
\end{proof}

\noindent
So the heuristic decides how much of the front has to be \emph{discovered}, and
nothing about what is returned.

\subsection{Which heuristic, and why the verification is free}

Two sources of efficient points were compared.

The \textbf{augmented weighted Tchebychev program} is the natural candidate:
its optima are efficient by construction, for any strictly positive weights and
any positive augmentation, so its output needs no verification.  A
\textbf{Pareto local search} with an archive proves nothing, so every member
has to be put through the exact efficiency test.

That verification looks like the deciding cost and is not, because of an
accounting that holds only here: \textbf{the enumeration was going to pay the
same efficiency test anyway}.  Every probe a seed removes was going to be
followed by one.  The verification cancels, and the probe is the gain.  What a
seed really costs is its share of the walk.

\begin{table}[h]
\centering
\caption{The two sources, per efficient point.}
\label{tab:sources}
\begin{tabular}{lrrr}
\toprule
 & cost of one point & verification & front covered \\
\midrule
Tchebychev scalarisation & $20.90$\,ms & none needed & $33$--$100\%$ \\
Pareto walk (whole archive) & $0.011$\,s total & one test each & $95\%$ median \\
\bottomrule
\end{tabular}
\end{table}

\noindent
And a scalarisation is not cheaper than the probe it replaces: one probe costs
$9.48$\,ms and $0$ branch \& bound nodes, one scalarisation $20.90$\,ms and
$21$ nodes, cheaper on only $7$ of $40$ instances.  The probe maximises a
linear direction over $\Dset$ plus ordinary rows and usually needs no branching
at all; the scalarisation minimises a \emph{maximum}, at the price of a
continuous variable and $p$ rows, and it branches.

Consistently with that arithmetic, seeding from the generator \emph{loses}, and
loses \emph{more the more of the front it seeds}: with a wide weight grid it
covers $100\%$ of the front and runs at $2.61\times$ the unseeded time.

\subsection{The order in which seeds are handed over}
\label{sec:order}

Placing $k$ seeds was expected to remove $k$ probes.  In the median it does
($+1.00$ over $152$ instances), but on $5$ of those $152$ the seeded run spent
\emph{more}: the pre-split imposes a box structure the search would not have
chosen, and a box it would have absorbed into a larger one now exists
separately and must be probed.

The probe maximises a fixed direction, so the search splits first around
whatever scores highest on it.  Sorting the seeds by that same direction
reproduces the structure the search would have built.

\begin{table}[h]
\centering
\caption{Probes against the unseeded enumeration, by the order the seeds
arrive.  Twelve instances, $n = 6 \dots 10$.}
\label{tab:order}
\begin{tabular}{lrr}
\toprule
order & probes & raised the count on \\
\midrule
archive order        & $0.89\times$ & $4$ of $12$ \\
\textbf{probe order} & $\mathbf{0.70\times}$ & \textbf{none} \\
reversed             & $0.94\times$ & $5$ of $12$ \\
\bottomrule
\end{tabular}
\end{table}

\noindent
The reversed arm is what makes the direction causal rather than a coincidence
of this sample.  In archive order the hybrid is \emph{slower} than the pure
exact method, at $0.81\times$; in probe order it is $1.33\times$ faster.

\section{Computational results}
\label{sec:results}

All arithmetic is exact and rational, with no floating point anywhere, so every
comparison is between two exactly solved problems.  Instances are drawn from a
documented pseudo-random family; the statistic is the \emph{median of the
per-instance ratios} with the interquartile range, never a ratio of totals,
which would be decided by the single heaviest instance.

\begin{table}[h]
\centering
\caption{The hybrid against the pure exact methods.  Thirty instances at each
of five configurations, $n = 6 \dots 9$, $p = 3,4$.}
\label{tab:hybrid}
\begin{tabular}{lrrr}
\toprule
 & median & IQR & faster on \\
\midrule
front enumeration        & $1.33\times$ & $[1.20,\ 1.51]$ & $139/150$ \\
$\E$ complete            & $\mathbf{1.31\times}$ & $[1.17,\ 1.47]$ & $139/150$ \\
probes                   & $0.72\times$ & $[0.67,\ 0.75]$ & \emph{never worse} \\
\bottomrule
\end{tabular}
\end{table}

\noindent
The answer is identical on every one of the $150$ --- same front, same $\E$,
same optimum of $(P_E)$ --- which the experiment asserts rather than assumes.

For context, the underlying substitution of criterion space for decision space
is worth the following on $(P_E)$ itself, thirty instances per row at $n=6$:

\begin{table}[h]
\centering
\caption{Decision space against criterion space on $(P_E)$.}
\label{tab:margin}
\begin{tabular}{crrrr}
\toprule
$p$ & ratio (median) & IQR & sub-programs & b\&b nodes \\
\midrule
2 & $1.52\times$  & $[1.25,\ 1.91]$  & $1.72\times$ & $1.28\times$ \\
3 & $2.69\times$  & $[1.87,\ 3.72]$  & $1.40\times$ & $1.79\times$ \\
4 & $4.99\times$  & $[2.39,\ 11.62]$ & $1.19\times$ & $2.68\times$ \\
5 & $7.93\times$  & $[3.61,\ 19.98]$ & $1.17\times$ & $3.16\times$ \\
6 & $14.57\times$ & $[7.90,\ 25.17]$ & $\mathbf{1.02\times}$ & $\mathbf{5.37\times}$ \\
\bottomrule
\end{tabular}
\end{table}

\noindent
The last row is the cost model of Section~\ref{sec:cost} read off the
measurement: at $p=6$ the two methods solve \emph{the same number} of
sub-programs and one still takes $14.57\times$ as long, so the whole margin
lies in the cost of a single sub-problem, whose branch \& bound explores
$5.37\times$ the nodes because it carries $p$ binaries and $p+1$ big-$M$ rows
per cut.

\paragraph{Measurement discipline.}  The machine used has a noise floor of
roughly $\pm 30\%$ per instance even at per-instance minima over five repeats,
so a small effect cannot be resolved in seconds and no number of repeats fixes
that.  Deterministic work counts --- sub-programs, branch \& bound nodes,
simplex pivots --- are reported next to the times and are identical on every
run and every machine.  Runs stopped by a time budget are reported as censored
and excluded from medians rather than counted as their budget.

\section{Negative results, and a wall}
\label{sec:wall}

We report these at the same length as the positive ones, because they delimit
where hybridisation can help at all.

\subsection{Four changes that remove a quarter of the sub-problems and buy
nothing}

\begin{table}[h]
\centering
\caption{Four attempts on the sub-problem count.}
\label{tab:wall}
\begin{tabular}{lrr}
\toprule
change & sub-problems removed & cost \\
\midrule
early exit on a hopeless bound     & $25\%$       & no measurable time \\
free contradiction filter          & $24.8\%$     & $0.98\times$ pivots \\
range filter (ideal, anti-ideal)   & $+11.3\%$    & $1.03\times$ pivots \\
dominated-corner filter            & $8$--$23\%$  & $1.077\times$ pivots \\
\bottomrule
\end{tabular}
\end{table}

\noindent
All four are exact and one-sided: they discard only boxes that provably hold
nothing, so they can change the cost and not the answer.  Only the second is
kept, and it is worth $2\%$.  After four attempts the reason is a property
rather than an accident:

\begin{quote}
\textbf{In this search, removing sub-problems does not remove time, because the
sub-problems that can be removed cheaply are the ones that were already cheap.}
\end{quote}

\noindent
$73\%$ of all boxes are empty, and an empty box is refuted by the simplex
almost at once: every probe in these families costs \emph{zero} branch \& bound
nodes, because the relaxation is integral or infeasible at the root.  A filter
that catches empty boxes is catching the cheapest work in the search.  The only
change reported here that goes round this wall is Section~\ref{sec:elim},
which attacks an \emph{exponent} rather than a count.

\subsection{A single-program completeness certificate}

After the walk seeds $95\%$ of the front, $93\%$ of the remaining boxes are
empty and the enumeration is almost purely a \emph{proof of emptiness}.  But
``is there a feasible point not dominated by any of the $k$ known ones?'' is
\textbf{one} question, and the box list answers it by splitting it into
hundreds.  Formulation \eqref{eq:cut} answers it in one integer program with
$k$ cuts, built once and never iterated --- which removes exactly the growth
that makes that formulation expensive in Section~\ref{sec:cost}.

It is sound: $13$ of $30$ instances were proved complete by a single program,
with no wrong verdict.  It is also $4.5\times$ slower in the median and
$100\times$ slower at worst ($239$\,s against $1.95$\,s), because $k$ cuts mean
$kp$ binaries in one model --- $88$ of them at $k = 22$, $p = 4$.

This answers a question wider than an optimisation.  Criterion space is the
better \emph{search} because the model never grows; the certificate tests
whether it is also the better \emph{proof}, where there is no iteration and no
growth at all.  \textbf{It is, by $4.5\times$}: answering many cheap questions
beats answering one question with $88$ binaries.

\subsection{What the failures have in common}

Four hybridisations were attempted in this line of work.  Handing the exact
search the unfinished work of the decision-space method ties it.  Seeding the
optimisation search from the Tchebychev program never pays, because that
program steers in criterion space relative to the ideal point and never looks
at $\Phi$.  Seeding the \emph{enumeration} from the same program loses, because
a scalarisation costs more than the probe it replaces.  Only the Pareto walk
pays, and only in probe order.  The rule that survives all four:

\begin{quote}
A hybridisation pays only where the exact method has a mechanism the first
phase can feed \textbf{more cheaply than the exact method feeds it itself}.
\end{quote}

\noindent
The first three fail the second half of that test.  The fourth passes it
because its verification is a test the exact method had already budgeted for.

\section{Conclusion}

Keeping the unexplored region in criterion space removes the growth that
dominates decision-space methods for $(P_E)$, and the same substitution carries
to complete enumeration of the non-dominated set.
Theorem~\ref{thm:slice} then turns that enumeration into the complete efficient
set at no efficiency test at all, and solving the slice equations rather than
pruning on them makes the remaining enumeration $49\times$ cheaper.  On top of
that, a Pareto local search whose verified archive is handed over in probe
order reaches all of $\E$ at $1.31\times$ the pure exact method over $150$
instances, with identical answers.

The negative results are the boundary of the positive one.  A seed is worth
what the exact method can do with it, and in an enumeration --- which has no
bound to feed --- it is worth exactly the discovery it replaces, priced against
what discovery would have cost.

\paragraph{Limitations.}  The instances are drawn from one pseudo-random family
with $n \leqslant 12$ and $p \leqslant 6$.  The \emph{directions} reported here
are explained by mechanisms that are themselves measured --- sub-program counts,
node counts, pivot counts --- but the \emph{magnitudes} are properties of that
family and should not be read as expected values elsewhere.  A full
bibliographic review, including sources behind paywalls, remains to be done
before any claim of novelty is made: what is stated above is what we measured,
not what we believe to be new.

\paragraph{Reproducibility.}  The implementation is pure Python with no
third-party dependency, computes in exact rational arithmetic throughout, and
carries $92$ tests.  Several check \emph{properties} rather than answers --- the
split is a partition, the metaheuristic never hands over an inefficient point,
a budgeted run never claims completeness --- since a method that silently lost
efficient points would still agree with the optimum on favourable instances.
Per-instance records of every experiment are stored so that each table can be
rebuilt and re-checked without re-running it.

\end{document}
